% ================================================================
% Section: The Public Evaluation Record
% Documents what the public reporting record can and cannot support,
% and why a cross-model accuracy table cannot be built from it.
% No original benchmark runs.
% ================================================================

\section{The Public Evaluation Record for Small Language Models}
\label{sec:reported}

This section examines what the public record can support about \gls{slm} capability and deployment
behavior. In contrast to an original benchmarking study,
the performance values analyzed here are \emph{reported} figures, synthesized
from vendor technical reports, official model cards, community leaderboards, and
third-party evaluations. We first document the protocol used to collect and
reconcile these figures (Section~\ref{subsec:reported_protocol}), then present
the consolidated synthesis (Sections~\ref{subsec:reported_overall}--\ref{subsec:reported_safety}).
Appendix~\ref{app:provenance} records the provenance of these figures at the
level of model and source class, together with the reference conditions under
which they were reported; where sources disagree, the reconciliation rule of
Section~\ref{subsec:reported_sources} applies.

\subsection{Data-Collection Protocol}
\label{subsec:reported_protocol}

\subsubsection{Model Selection}
\label{subsec:reported_model_selection}
Model inclusion required: (1)~public availability of weights or a publicly
documented evaluation card; (2)~parameter count below 10~billion, with Phi-4 at
14B retained as an upper-bound reference; (3)~at least one reported score on an
established benchmark used here; and (4)~representation of diverse training
methodologies (synthetic-data-centric, web-scale pretraining, distillation).
Table~\ref{tab:model_specs} lists the resulting 18-model cohort with
architectural specifications as reported by each model's release materials. The
cohort spans from 0.5B to 14B parameters and eight developer organizations,
enabling analysis of how training methodology, architecture, and scale interact
to determine reported capability.

\begin{table}[ht]
\centering
\caption{Architectural specifications of reviewed models (as reported by each model's official release materials).}
\label{tab:model_specs}
\footnotesize\adjustbox{max width=\linewidth}{%
\begin{tabular}{lccccccc}
\toprule
\textbf{Model} & \textbf{Params} & \textbf{Layers} & \textbf{Hidden} & \textbf{Attention} & \textbf{Context} & \textbf{Vocab} & \textbf{Train Tokens} \\
\midrule
TinyLlama-1.1B & 1.1B & 22 & 2,048 & MHA-32h & 2K & 32K & 3.0T \\
SmolLM2-1.7B & 1.7B & 24 & 2,048 & MHA-32h & 2K & 49K & 11.0T \\
StableLM 2-1.6B & 1.6B & 24 & 2,048 & MHA-32h & 4K & 100K & 2.0T \\
Qwen2.5-0.5B & 0.5B & 24 & 896 & GQA-14/2 & 32K & 152K & 18.0T \\
Qwen2.5-1.5B & 1.5B & 28 & 1,536 & GQA-12/2 & 32K & 152K & 18.0T \\
Qwen2.5-3B & 3.0B & 36 & 2,048 & GQA-16/2 & 32K & 152K & 18.0T \\
Qwen2.5-7B & 7.6B & 28 & 3,584 & GQA-28/4 & 128K & 152K & 18.0T \\
Llama 3.2-1B & 1.2B & 16 & 2,048 & GQA-32/8 & 128K & 128K & 9.0T \\
Llama 3.2-3B & 3.2B & 28 & 3,072 & GQA-24/8 & 128K & 128K & 9.0T \\
Gemma 2-2B & 2.6B & 26 & 2,304 & GQA-8/4 & 8K & 256K & 2.0T \\
Gemma 2-9B & 9.2B & 42 & 3,584 & GQA-16/8 & 8K & 256K & 8.0T \\
Gemma 3-1B & 1.0B & 26 & 1,536 & GQA-LG & 32K & 262K & 2.0T \\
Gemma 3-4B & 4.3B & 34 & 2,560 & GQA-LG & 128K & 262K & 4.0T \\
Mistral 7B & 7.3B & 32 & 4,096 & GQA-32/8 & 8K & 32K & n/a \\
Phi-3-mini & 3.8B & 32 & 3,072 & MHA-32h & 128K & 32K & 3.3T \\
Phi-3.5-mini & 3.8B & 32 & 3,072 & MHA-32h & 128K & 32K & 3.4T \\
Phi-4-mini & 3.8B & 32 & 3,072 & GQA-24/8 & 128K & 200K & 5.0T \\
Phi-4 & 14.0B & 40 & 5,120 & GQA & 16K & 100K & $\sim$10T \\
\bottomrule
\end{tabular}%
}\par\vskip 4pt\relax
{\raggedright \small \textit{Note:} MHA = Multi-Head Attention; \gls{gqa} = Grouped-Query Attention (Q/KV heads); GQA-LG = \gls{gqa} with local--global interleaving. Every model in this cohort is dense, so total and active parameters coincide and a single parameter column suffices; no small-active-parameter \gls{moe} model (Section~\ref{subsubsec:intro_def_terms}) is in the cohort, because none has public reporting complete enough to include.

\textit{On training-token counts.} These are the least reliable column in the
table, and their status differs by row. Marked \emph{per-variant} where the
source states a figure for that specific size: Gemma~2-2B (2.0T), Gemma~2-9B
(8.0T), Gemma~3-1B (2.0T), Gemma~3-4B (4.0T), Phi-3-mini (3.3T), Phi-3.5-mini
(3.4T), Phi-4-mini (5.0T), TinyLlama (3.0T), SmolLM2-1.7B (11.0T), StableLM~2
(2.0T). Marked \emph{series-level} where the source reports one figure for the
whole family and the individual sizes are not separately stated: Qwen2.5
(18.0T across 0.5B--7B) and Llama~3.2 (9.0T across 1B and 3B). Phi-4's
${\sim}10$T is an approximation given as such by its source. Mistral~7B's
training-token count is not disclosed. We report series-level figures in the
size rows they cover rather than omitting them, but they should not be read as
independent per-model measurements, and no claim in this paper relies on
comparing a series-level figure against a per-variant one.

\textit{On context length.} The value is the maximum the primary source
documents, which for several models is reached through position interpolation at
inference rather than trained natively: Qwen2.5-7B is natively 32K and extended
to 128K, and Phi-3-mini's 128K is a separate long-context checkpoint from the 4K
base.\par}
\end{table}

\subsubsection{Benchmark Tasks and Metrics}
\label{subsec:reported_benchmarks}
The benchmarks that recur across the reviewed release documents span seven
capability areas: \gls{mmlu} \citep{hendrycks2021mmlu} for broad knowledge;
ARC-Easy and ARC-Challenge \citep{clark2018arc} and HellaSwag
\citep{zellers2019hellaswag} for commonsense and language reasoning; \gls{gsm8k}
\citep{cobbe2021gsm8k} and MATH \citep{hendrycks2021math} for mathematical
reasoning; HumanEval \citep{chen2021humaneval} and MBPP \citep{austin2021mbpp}
for code generation; TriviaQA \citep{joshi2017triviaqa} and SQuAD~2.0
\citep{rajpurkar2018squad2} for question answering; MNLI \citep{williams2018mnli}
for natural language inference; and SST-2 \citep{socher2013sst} for sentiment
classification. Each benchmark carries a standard metric: accuracy for multiple
choice, exact match or F1 for question answering, pass@1 for code, exact match
for \gls{gsm8k}, and symbolic equivalence for MATH. Instruction-following
(MT-Bench-style \citep{zheng2023mtbench}) and summarization are reported too
sparsely and too heterogeneously to compare at all, and the conventions behind
them are described in Appendix~\ref{subsec:appendix_human_eval}. Section~\ref{subsec:reported_agentic_gap}
takes up separately the question of whether any of these benchmarks bears on
agentic suitability.

Four conditions determine what a benchmark number means, and a source that omits
any of them has not reported a comparable result. \emph{Benchmark version:}
\gls{mmlu} exists in the original and in MMLU-Pro, ARC in Easy and Challenge
splits, GSM8K in original and GSM8K-Platinum revisions, and HumanEval in the
original and HumanEval+ forms, and these are not interchangeable.
\emph{Prompt template:} whether the model's own chat template is applied,
whether options are lettered or restated, and whether a system prompt is present.
\emph{Answer extraction:} whether a multiple-choice answer is scored by
comparing per-option log-likelihoods or by parsing a generated letter, which can
differ by several points on the same model and benchmark, and for
chain-of-thought tasks, which span of the generation is taken as the answer.
\emph{Decoding settings:} greedy versus sampled, temperature, and for pass@$k$
the number of samples. \emph{Harness:} the framework and its version, most often
\texttt{lm-evaluation-harness} \citep{gao2023harness} or HELM
\citep{liang2023helm}, each with its own defaults for all of the above.

The coverage of these five conditions across the reviewed sources is what
Section~\ref{subsec:reported_overall} reports on, and it is poor. Which
benchmarks a vendor reports, in which version, under how many shots, and with
which extraction rule varies enough to prevent a common table.

\subsubsection{Hardware and Quantization Scope}
\label{subsec:reported_hardware}
The review covers reported behavior across six configurations spanning the
cloud-to-edge spectrum (Table~\ref{tab:hardware_configs}) and three precisions:
FP16 (reference), INT8, and INT4, with \gls{gguf}~Q4\_K\_M, GPTQ group-128, and \gls{awq}
group-128 treated as distinct INT4 subclasses when the source distinguishes them.
Inference-framework attribution (vLLM \citep{kwon2023vllm}, llama.cpp
\citep{gerganov2023llamacpp}, ONNX Runtime \citep{microsoft2019onnx}, ExecuTorch
\citep{meta2024executorch}, MLC-LLM \citep{mlc2024}) follows the stack used by the
source of each figure; cross-framework comparisons are made only when compared
figures share a framework or are explicitly flagged.

\begin{table}[ht]
\centering
\caption{Hardware configurations covered by the review. For each reviewed model, reported performance on a configuration was included where publicly available.}
\label{tab:hardware_configs}
\footnotesize\adjustbox{max width=\linewidth}{%
\begin{tabular}{lllll}
\toprule
\textbf{Configuration} & \textbf{Hardware} & \textbf{RAM/VRAM} & \textbf{TDP} & \textbf{Use Case} \\
\midrule
Cloud GPU & NVIDIA A100 80GB & 80 GB HBM2e & 300W & Training + serving \\
High-end Consumer & NVIDIA RTX 4090 & 24 GB GDDR6X & 450W & Local inference \\
CPU/iGPU Laptop & AMD Ryzen 7 7840HS & 32 GB DDR5 & 35--54W & Laptop deployment \\
Edge GPU Module & NVIDIA Jetson Orin Nano & 8 GB LPDDR5 & 25W & Edge AI / robotics \\
Edge SBC & Raspberry Pi 5 & 8 GB LPDDR4X & 12W & Embedded IoT \\
Mobile SoC & Snapdragon 8 Gen 3 & 12 GB LPDDR5X & $\sim$8W sustained & Smartphone \\
\bottomrule
\end{tabular}%
}
\end{table}

\subsubsection{Source Selection, Extraction, and Reconciliation}
\label{subsec:reported_sources}
Reported values were collected in decreasing order of priority from: (1)~official
vendor technical reports and release papers; (2)~official model cards on Hugging
Face and vendor repositories, including linked evaluation appendices;
(3)~community-operated standardized evaluations (the Open \gls{llm} Leaderboard, HELM
\citep{liang2023helm}, and lm-evaluation-harness \citep{gao2023harness}-based
community reports); (4)~third-party independent evaluations (Artificial Analysis,
LMSYS Chatbot Arena, peer-reviewed follow-up studies); and (5)~vendor developer
blogs, used only to fill gaps. All figures were extracted within a single
collection window to preserve time-consistency across models.

When qualifying sources disagreed on the same model--task--condition tuple, the
figure from the source higher in the priority list was retained and the
alternative recorded in Appendix~\ref{app:provenance}; where the discrepancy
exceeded 2.0~percentage points (or the metric equivalent), both figures were
retained as a reported range or the tuple was excluded. A tuple was excluded from
cross-model comparison when: (a)~only a single self-reported value existed and
could not be cross-checked; (b)~the harness, prompt format, or shot count
materially differed from the setting used by at least half the compared models
and could not be normalized; or (c)~the source did not specify base versus
instruction-tuned status.

\subsubsection{Evidence Grading}
\label{subsec:reported_grading}

Because this survey mixes controlled studies, vendor self-reports, independent
reproductions, and our own synthesis, and because those carry very different
weight, every substantive quantitative claim in the paper is graded. The scheme
has four levels, and the grade denotes the strongest class of source standing
behind the claim, not the claim's importance:

\begin{description}
  \item[\evA{} Controlled] Measured under a single named harness, with the
    compared conditions differing only in the variable of interest, and a matched
    baseline reported. This is the standard a cross-model comparison would have
    to meet. \emph{No claim in this paper carries grade~\evA}, which is the
    principal finding of Section~\ref{subsec:reported_overall}.
  \item[\evB{} Independent] Reported by a party with no stake in the result:
    a peer-reviewed third-party study, a standardized community evaluation, or a
    reproduction by an unaffiliated group.
  \item[\evC{} Reported] Stated by the developer of the artifact or the author of
    the technique, in a technical report, model card, or release paper, without
    independent reproduction.
  \item[\evD{} Interpretation] Our synthesis across sources, or an inference from
    evidence that was not gathered to answer the question being asked. This grade
    is applied wherever we combine figures across sources, and it is the correct
    grade for most directional statements in this survey.
\end{description}

Grades appear in the tables that carry quantitative claims
(Tables~\ref{tab:compression_quality} and~\ref{tab:agentic_evidence}) and against
the main claims of the paper in Table~\ref{tab:claim_grades}. Two consequences
are worth stating plainly. First, a grade~\evC{} claim is not thereby false; it
is unverified, and vendor self-reports in this field are systematically more
favorable than third-party reproductions, so the direction of the likely error is
known. Second, the absence of any grade~\evA{} claim is not a defect peculiar to
this survey but a property of the field's current evidence base, and it is the
reason Section~\ref{subsec:challenges_verification} is written as a specification
for experiments rather than a list of topics.

\begin{table}[t]
\centering
\caption{Evidence grade for the principal claims of this survey. Grades follow Section~\ref{subsec:reported_grading}: \evA~controlled, \evB~independent, \evC~reported by the interested party, \evD~our interpretation.}
\label{tab:claim_grades}
\footnotesize\adjustbox{max width=\linewidth}{%
\begin{tabular}{@{}p{9.2cm}cl@{}}
\toprule
\textbf{Claim} & \textbf{Grade} & \textbf{Where} \\
\midrule
A cross-model accuracy table cannot be assembled from current public reporting without misrepresenting its sources & \evB & \S\ref{subsec:reported_overall} \\
The strongest 3--4B models of 2026 report scores that were state of the art for 13B-class models in 2023 & \evC & \S\ref{subsec:reported_overall} \\
Within families that publish ablations, corpus composition moves results more than architecture at fixed scale & \evD & \S\ref{subsec:discussion_key_findings} \\
INT8 post-training quantization is effectively lossless on knowledge and mathematical benchmarks & \evB & \S\ref{subsec:reported_quantization} \\
INT4 accuracy loss grows as model size falls and is larger for multi-step reasoning & \evC & \S\ref{subsec:reported_quantization} \\
Quantization preserves fidelity better than pruning below 4B parameters at matched compression & \evB & \S\ref{subsec:compression_quantization} \\
Weight footprint at 4~bits is $3.2$--$3.4\times$ below FP16 & Derived & \S\ref{subsec:reported_quantization} \\
Sub-10B function-calling models are competitive with far larger models on function-calling benchmarks & \evB & \S\ref{subsec:agentic_evidence} \\
Serving a 7B model is 10--30$\times$ cheaper than serving a 70--175B model & \evC & \S\ref{subsubsec:agentic_economics} \\
End-to-end agent success rates for sub-10B models on realistic multi-step benchmarks are low in absolute terms & \evB & \S\ref{subsec:agentic_evidence} \\
Hypothesis~H: a specialized SLM is the appropriate default for bounded, recurrent, short-context, verifiable agent calls & \evD & \S\ref{subsec:intro_thesis} \\
\bottomrule
\end{tabular}%
}
\end{table}

\subsubsection{Limitations of the Evidence Base}
\label{subsec:reported_limits}
Because the figures are compiled from heterogeneous public sources rather than a
single controlled pipeline, four limitations bound the strength of every
comparison below. \textbf{Prompt-format and harness variation:} different sources
use different templates and answer-extraction heuristics, which can shift reported
accuracy by several points under the same nominal benchmark. \textbf{Self-reporting
bias:} vendor-reported scores are systematically more favorable than third-party
reproductions; where third-party figures exist, comparisons preferentially use
them and flag the divergence. \textbf{Condition-coverage sparsity:} not every
model has reported scores for every task, tier, or precision. Missing values are
never imputed, interpolated, or carried across from a sibling model size or from
a base checkpoint to its instruction-tuned variant; a tuple with no qualifying
source is recorded as absent, and absence is itself reported, since the pattern
of what vendors decline to publish is part of the finding of
Section~\ref{subsec:reported_overall}. No statistic in this paper is computed
over a set from which absent tuples were silently dropped. \textbf{Snapshot staleness:} the ecosystem evolves monthly,
and figures reflect the state of public reporting at the collection window. These
limitations motivate the empirical-verification agenda of
Section~\ref{subsec:challenges_verification}.

\subsection{Why a Common Benchmark Table Cannot Be Built from Public Reporting}
\label{subsec:reported_overall}

Our original intent for this section was a single consolidated accuracy table
placing all eighteen models side by side on a shared benchmark set. Assembling it
from the public record proved impossible without misrepresenting the sources, and
the reasons are systematic rather than incidental. Because that finding bears
directly on how practitioners should read vendor claims, we report it in place of
the table.

\emph{Vendors do not report a common benchmark set.} Microsoft reports
\gls{mmlu}, \gls{gsm8k}, ARC-Challenge, HellaSwag and MATH for Phi-4-mini
\citep{microsoft2025phi4mini}. Alibaba's instruction-tuned Qwen2.5 evaluations
report MMLU-Pro, \gls{gsm8k}, MATH and HumanEval, but not ARC-Challenge or
HellaSwag \citep{qwen2024qwen25}. Google reports ARC-c, HellaSwag, \gls{mmlu},
\gls{gsm8k}, MATH, HumanEval and MBPP for Gemma, but its headline tables cover
the pre-trained checkpoints \citep{gemma2team2024,team2025gemma3}. The
intersection across any six vendors is close to empty, so a full matrix can only
be completed by mixing sources of different provenance into the same row.

\emph{Shot counts differ per benchmark and per vendor.} Within a single Phi-4-mini
model card, \gls{mmlu} is 5-shot, \gls{gsm8k} is 8-shot with chain-of-thought,
ARC-Challenge is 10-shot and HellaSwag is 5-shot. Gemma reports ARC-c at 25-shot,
HellaSwag at 10-shot, MATH at 4-shot and MBPP at 3-shot. Llama~3.2 reports
ARC-Challenge and HellaSwag zero-shot but \gls{mmlu} at 5-shot. A column headed
with a single evaluation protocol is therefore false for most of its cells, and
the differences are large: shot count alone moves reported accuracy by several
points on multiple-choice tasks.

\emph{Base and instruction-tuned figures are routinely conflated.} Several
families publish their most complete benchmark tables for the pre-trained
checkpoint while shipping the instruction-tuned variant as the artifact
practitioners actually deploy. Gemma~3-4B is reported at 59.6 \gls{mmlu} and 38.4
\gls{gsm8k} after pre-training \citep{team2025gemma3}; the instruction-tuned model
is described qualitatively as far stronger, without an equivalently detailed
table. Any row that silently mixes the two overstates the weaker model and
understates the stronger one.

\emph{Matched baselines are rarely published.} Quantized, fine-tuned and
long-context variants are seldom reported alongside an FP16 baseline evaluated
under the same harness, so the deltas that matter most for deployment decisions
cannot be recovered from the sources at all.

The consequence is that the widely circulated cross-model \gls{slm} comparison
tables, including those in secondary sources and community leaderboards, should
be treated as approximate at best. Where a genuine comparison is needed, the only
sound basis is a controlled re-run under one harness, which is why
Section~\ref{subsec:challenges_verification} sets that out as the first item of
the verification agenda rather than as an optional extension. What can be stated
from the public record without distortion is qualitative and still substantial:
the strongest 3--4B models of 2026 report scores on mathematics, code and
knowledge benchmarks that were state of the art for 13B-class systems in 2023;
no single \gls{slm} leads across all task families; and the largest single source
of variation among models of near-identical size is the training corpus rather
than the architecture, a pattern visible in every family's own ablations even
though it cannot be quantified across families.

\subsection{Why These Benchmarks Do Not Answer the Agentic Question}
\label{subsec:reported_agentic_gap}

There is a second problem with the record, independent of its internal
inconsistency, and it bears directly on the thesis of this paper. Every benchmark
in Section~\ref{subsec:reported_benchmarks} measures a model answering a question
in isolation. None measures what an agent actually asks a model to do.

\gls{mmlu} measures recall of facts under a multiple-choice format. \gls{gsm8k}
measures arithmetic word problems solved in one turn. HumanEval measures whether
a standalone function passes unit tests. An agent, by contrast, asks the model to
select one of several available tools, to emit arguments that satisfy a schema,
to remain well-formed across a long tail of inputs, to recover when a tool
returns an error, and to do so repeatedly without a single failure derailing the
trajectory. A model can be strong on all three of the benchmarks above and still
emit malformed JSON often enough to be unusable in a loop, and the reverse holds
too: a small model tuned for function calling can be weak on \gls{mmlu} and
entirely adequate as an agent's executor.

The consequence is that the entire evaluation record examined in this section,
even if it were internally consistent, would be the wrong evidence for
hypothesis~H. Section~\ref{subsec:agentic_evidence} therefore assembles a
separate body of agent-specific evidence, drawn from function-calling, tool-use,
structured-output, and end-to-end agent benchmarks, and grades it under the same
scheme. We flag here that this second body of evidence has its own gaps, which
are different from and in some respects worse than the ones documented above:
the agent benchmarks that exist report frontier models far more often than
sub-10B ones, and almost none reports cost per completed task, which is the
quantity the economic half of the thesis turns on.

\subsection{Quantization Impact}
\label{subsec:reported_quantization}

Across the compression literature reviewed in Section~\ref{sec:compression}, the
reported picture for the sub-10B regime is consistent on four points. INT8
post-training quantization is effectively lossless on knowledge and mathematical
benchmarks. INT4 costs a small but real amount of accuracy, and the loss is
\emph{not} uniform across model sizes: smaller models degrade more, consistent
with the argument that each weight in a small model carries proportionally more
information \citep{slmquant2025}. Mathematical and multi-step reasoning degrade
faster than knowledge recall under INT4, which matters directly for agentic use,
since arithmetic and tool-argument construction are exactly the operations agents
delegate. Among INT4 methods, \gls{awq} retains more accuracy than GPTQ by
protecting the small fraction of activation-salient channels
\citep{lin2024awq}, while \gls{gguf}'s Q4\_K\_M mixed-precision scheme gives the
best size-quality balance on CPU and edge backends
\citep{gerganov2023llamacpp}. Across the sub-4B range, quantization preserves
fidelity better than pruning at matched compression
\citep{kurz2025pruning_vs_quant}.

Per-model, per-benchmark degradation figures are not reported consistently
enough across this cohort to tabulate: sources differ in harness, calibration
set, and group size, and few publish a matched FP16 baseline alongside the
quantized result. Measuring them under a single toolchain is one of the
verification studies proposed in
Section~\ref{subsec:challenges_verification}.

The memory consequence, by contrast, is firmly established, because it follows
from the storage format rather than from measurement.
Table~\ref{tab:memory_footprint} gives weight footprints computed from parameter
count and bit-width. The realized reduction falls short of the nominal
4$\times$, at roughly 3.2 to 3.4$\times$, because embedding, normalization, and
in mixed schemes a subset of attention tensors are retained at higher precision.
A 7B model drops from about 15~GB to about 4.5~GB, and a 3.8B model from 7.6~GB
to about 2.4~GB, which is what moves each of them across a hardware tier
boundary.


\begin{table}[ht]
\centering
\caption{Model-weight memory footprint by precision (GB), computed from parameter count and storage bit-width. Four-bit columns include the higher-precision tensors each scheme retains, which is why the reduction is below the nominal 4$\times$. Activation and KV-cache memory are additional and depend on context length and batch size.}
\label{tab:memory_footprint}
\footnotesize
\begin{tabular}{lcccc}
\toprule
\textbf{Model} & \textbf{FP16} & \textbf{GPTQ-4bit} & \textbf{AWQ-4bit} & \textbf{Q4\_K\_M} \\
\midrule
Qwen2.5-7B & 15.2 & 4.2 & 4.4 & 4.5 \\
Phi-4-mini & 7.6 & 2.2 & 2.3 & 2.4 \\
Llama 3.2-3B & 6.4 & 1.9 & 2.0 & 2.1 \\
SmolLM2-1.7B & 3.4 & 1.1 & 1.1 & 1.2 \\
TinyLlama-1.1B & 2.2 & 0.7 & 0.8 & 0.8 \\
\bottomrule
\end{tabular}
\end{table}

\subsection{Multilingual and Safety Reporting}
\label{subsec:reported_safety}

Multilingual and safety reporting is markedly sparser than benchmark reporting
across the cohort, and the coverage that exists is not comparable across models;
we therefore report the direction of the evidence rather than tabulating it.
On multilingual capability, the models whose technical reports document
large non-English pretraining mixtures, namely Qwen2.5 \citep{qwen2024qwen25}, Gemma~3
\citep{team2025gemma3}, and Aya \citep{ustun2024aya}, report substantially
stronger non-English performance than the English-centric families, and the gap
is widest on non-English mathematical reasoning, where both language coverage and
reasoning capacity bind at once. On safety, per-model toxicity evaluations are
reported inconsistently and are frequently absent altogether from research-grade
small models, so the cohort cannot be ranked; what the literature does establish
is that alignment is harder to preserve at reduced scale and that fine-tuning can
erode what alignment exists \citep{qi2023finetuning}, and that compression
degrades performance unevenly across languages and dialects
\citep{blodgett2020language,kurz2025pruning_vs_quant}. The absence of
comparable safety reporting is itself a finding, and it is one of the
verification gaps named in Section~\ref{subsec:challenges_verification}.
Both observations carry directly into
Section~\ref{sec:agentic}: an agent that routes non-English or safety-sensitive
calls to a small English-centric model inherits these gaps, which is an argument
for capability-aware rather than size-only routing.
